

\section{Le Grec Moderne comme ressource : Continuité lexicale et motivation}
space{0.3cm}
oindent

La présente étude s'inscrit dans le cadre de la section III.4.\cite{ref2} de la thèse de doctorat
portant sur la refondation cognitive de l'enseignement du grec ancien. Elle vise à examiner, de
manière exhaustive et critique, l'hypothèse selon laquelle le grec moderne et le fonds lexical
scientifique international ne constituent pas de simples adjuvants culturels, mais des ressources
cognitives fondamentales pour l'acquisition du grec ancien.

L'enseignement des langues anciennes traverse une crise de légitimité et d'efficacité, souvent
attribuée à des méthodes focalisées sur la grammaire-traduction qui isolent la langue cible de tout
contexte vivant. Cette isolationnisme pédagogique, qualifié par certains de « nécrophilie
linguistique », a historiquement coupé le grec ancien de ses deux ancrages naturels : sa descendance
vivante (le grec moderne) et sa diffusion universelle (le vocabulaire scientifique). Or, les
recherches récentes en linguistique historique, en psycholinguistique et en didactique des langues
suggèrent qu'un changement de paradigme est nécessaire. Ce rapport explore comment la réactivation
de la « continuité lexicale » permet de réduire la charge cognitive de l'apprentissage, de restaurer
la motivation des élèves par le sens, et de reconnecter l'objet scolaire à une réalité sociale et
culturelle tangible.

Nous articulerons notre analyse autour de quatre axes majeurs : la validation scientifique de la
continuité linguistique du grec à travers les âges ; l'analyse psycholinguistique du transfert
lexical et de l'hypothèse du seuil de lecture ; l'examen des méthodologies actives qui exploitent
cette continuité (notamment les approches de l'Institut Polis et d'Elliniki Agogi) ; et enfin,
l'exploitation didactique du vocabulaire scientifique comme levier de motivation intrinsèque dans le
cadre des programmes actuels.



\subsection{Épistémologie de la continuité linguistique : déconstruire la rupture}

L'obstacle premier à l'utilisation du grec moderne comme ressource est d'ordre épistémologique : la
perception tenace d'une rupture radicale entre le grec \og classique\fg{} (identifié quasi
exclusivement à l'attique des Ve et IVe siècles av. J.-C.) et les états ultérieurs de la langue.
Pour fonder une didactique de la continuité, il convient d'abord d'établir la réalité scientifique
de l'unicité de la langue grecque.



\subsection{L'unité diachronique du grec : perspectives historiographiques}

La linguistique historique contemporaine a largement remis en cause la vision segmentée de
l'histoire du grec. Les travaux fondateurs de Geoffrey Horrocks 1 et de Robert Browning 4
constituent le socle théorique indispensable à cette refondation.

Geoffrey Horrocks, dans son ouvrage de référence \textit{Greek: A History of the Language and its
Speakers}, démontre que le développement de la langue grecque, depuis la période mycénienne jusqu'au
grec moderne standard, forme un continuum ininterrompu. Il qualifie l'attique classique et le grec
moderne de « serre-livres » (\textit{book-ends}) d'une même histoire évolutive, plutôt que de
langues distinctes.\cite{ref1} Cette perspective est cruciale : elle postule que les changements
observés sont des évolutions internes au système et non des substitutions de systèmes. Francisco
Rodríguez Adrados renchérit en soulignant que le grec, avec le chinois, est l'une des deux seules
langues de culture connues depuis plus de 3500 ans qui soient encore parlées aujourd'hui.\cite{ref7}

Cette continuité n'est pas seulement une curiosité historique, elle a des implications didactiques
directes. Browning note que pour les locuteurs du grec moderne, les poèmes homériques ne sont pas
écrits dans une langue étrangère, mais dans une forme archaïque de leur propre langue.\cite{ref5}
Cette familiarité, même si elle nécessite une médiation scolaire, indique que la distance
linguistique est franchissable par des mécanismes de reconnaissance lexicale et morphologique,
contrairement à la distance qui sépare le français du latin, devenue infranchissable sans
apprentissage exogène complet.



\subsection{La question de la diglossie et la fossilisation}

L'un des facteurs ayant masqué cette continuité aux yeux des pédagogues occidentaux est le phénomène
de la diglossie. Comme l'analyse Horrocks, l'émergence précoce d'un canon classique (dès l'époque
hellénistique et romaine avec le mouvement atticiste) a eu un effet fossilisant sur la langue
écrite.\cite{ref3} Cette tension entre une langue parlée évolutive (\textit{Koine}, puis grec
médiéval et moderne) et une norme écrite archaïsante a créé une illusion de rupture.

Cependant, cette diglossie a paradoxalement servi la continuité lexicale. La \textit{Katharevousa}
(langue puriste du XIXe et XXe siècle), en réinjectant massivement du vocabulaire et de la
morphologie ancienne dans l'usage officiel, a maintenu vivants des termes et des structures qui
auraient pu disparaître naturellement.\cite{ref2} Ainsi, le grec moderne standard
(\textit{Dimotiki}) est aujourd'hui une fusion qui intègre un vaste héritage lexical ancien. Pour
l'enseignant, cela signifie que le vocabulaire \og moderne\fg{} est souvent, en réalité, du
vocabulaire ancien réactivé ou préservé. Ignorer cette ressource revient à se priver d'un accès
direct au lexique.



\subsection{Quantification du recouvrement lexical (overlap analysis)}

Au-delà des affirmations qualitatives, la refondation cognitive doit s'appuyer sur des données
probantes. Quel est le degré exact de similitude entre le grec ancien et le grec moderne? Les
estimations varient selon les corpus (Homère vs Nouveau Testament), mais convergent vers un taux de
recouvrement massif.

Plusieurs sources académiques et pédagogiques situent le recouvrement lexical entre le grec ancien
(Koine/Attique) et le grec moderne entre \textbf{70\% et 80\%}.\cite{ref11} Une analyse plus fine
est fournie par Jerneja Kavčič dans son étude sur la représentation du grec moderne dans les manuels
de grec ancien.\cite{ref15} En comparant le vocabulaire des manuels scolaires allemands, italiens et
anglophones avec le \textit{Dictionnaire du Grec Moderne Standard} (LKN), Kavčič établit une
typologie précise de la continuité :

\begin{table}[h!]\small\centering\begin{tabular}{| p{2.2cm} | p{2.8cm} | p{1.8cm} | p{3.2cm} |}\hline\textbf{Classe} & \textbf{Description} & \textbf{\%} & \textbf{Implication Cognitive} \\ \hline\textbf{Classe 1 : Identité} & Même forme et sens (\textit{thalassa}, \textit{anthropos}) & 50\% & Transfert positif direct \\ \hline\textbf{Classe 2 : Continuité} & Variation mineure (\textit{pater}/\textit{pateras}) & 12-16\% & Transfert avec ajustement \\ \hline\textbf{Classe 3 : Divergence} & Même forme, sens différent & 5-11\% & Interférence, vigilance requise \\ \hline\textbf{Classe 4 : Rupture} & Mots disparus/remplacés & 20-30\% & Apprentissage ex nihilo \\ \hline\end{tabular}\end{table}Ces données 15 sont fondamentales. Elles révèlent que plus de \textbf{60\%} du vocabulaire enseigné
dans les manuels de grec ancien est immédiatement accessible ou déductible pour un apprenant
maîtrisant le grec moderne. La continuité n'est donc pas une vague ressemblance, mais une structure
dominante du lexique.



\subsection{Psycholinguistique de l'apprentissage : charge cognitive et motivation}

La refondation cognitive de l'enseignement postule que l'échec des méthodes traditionnelles provient
d'une surcharge cognitive (trop de règles, trop de mots inconnus simultanément) et d'un déficit de
motivation. L'utilisation du grec moderne intervient ici comme un outil de régulation cognitive.



\subsection{L'hypothèse du seuil lexical (laufer \& nation) et la lecture extensive}

L'objectif ultime de l'enseignement des langues anciennes est la lecture fluide des textes. Or, les
travaux en psycholinguistique de Batia Laufer et Paul Nation ont établi une règle d'or, connue sous
le nom de \textbf{\og Threshold Hypothesis\fg{}} (Hypothèse du Seuil).\cite{ref16} Selon cette
théorie, validée empiriquement :

\begin{itemize}\item Un lecteur doit connaître \textbf{95\%} des mots d'un texte pour en comprendre le sens global sans aide extérieure.\item Il doit en connaître \textbf{98\%} pour une lecture de plaisir et une compréhension approfondie.\end{itemize}Dans l'enseignement traditionnel du grec ancien, l'élève, même après plusieurs années, atteint
rarement ce seuil sur des textes authentiques, ce qui le condamne au décryptage laborieux (le \og
mot à mot\fg{}) et empêche l'automatisation de la lecture.

L'apport du grec moderne est ici décisif. Comme le soulignent les discussions pédagogiques 12, la
maîtrise préalable ou concomitante du vocabulaire grec moderne permet d'atteindre ce seuil beaucoup
plus rapidement, particulièrement pour les textes de la Koine (Nouveau Testament, Septante,
historiens post-classiques). Le \og bruit\fg{} lexical (les mots inconnus) est drastiquement réduit
par la connaissance du fonds moderne, permettant au cerveau de se concentrer sur la syntaxe et les
mots rares.



\subsection{La propaédeutique par le moderne : l'approche "backdoor"}

Cette réalité a conduit à l'émergence d'une stratégie dite de la \og porte dérobée\fg{}
(\textit{Backdoor Approach}) : apprendre le grec moderne pour faciliter l'accès à
l'ancien.\cite{ref21} Bien que débattue, cette approche présente des avantages cognitifs majeurs :

1. \textbf{Immersion vivante :} Le grec moderne peut être appris par immersion, écoute, et
interaction sociale, ce qui est neurologiquement plus efficace que l'apprentissage par règles
explicites.

2. \textbf{Intuition grammaticale :} Même si la syntaxe a évolué (perte de l'infinitif, réduction
des cas), la structure fondamentale de la phrase, le système aspectuel (aoriste/présent/parfait) et
la logique des prépositions restent largement partagés.\cite{ref24} L'élève \og sent\fg{} la langue
grecque avant de l'analyser.

3. \textbf{Ressources disponibles :} L'abondance de matériaux pédagogiques modernes, de médias et de
locuteurs natifs offre un environnement d'apprentissage riche (\og input\fg{}) que le grec ancien
seul ne peut offrir.



\subsection{La problématique de la prononciation : érasme contre la réalité sonore}

Un obstacle majeur à l'utilisation de cette continuité est la barrière phonétique artificielle
érigée par la prononciation érasmienne. Instaurée en Occident à la Renaissance 26, cette
prononciation reconstituée (distinguant \textit{beta} \[b\] de \textit{vita} \[v\], \textit{eta}
\[ê\] de \textit{ita} \[i\]) coupe l'apprenant de la réalité sonore de la Grèce
actuelle.\cite{ref28}

Christophe Rico, doyen de l'Institut Polis, et d'autres réformateurs 29 plaident pour l'adoption
d'une prononciation historique (proche de la Koine ou du byzantin, donc très similaire au moderne)
pour plusieurs raisons cognitives :

\begin{itemize}\item \textbf{La boucle audio-phonatoire :} Pour mémoriser efficacement, le cerveau a besoin de lier le graphème au phonème de manière stable et vivante. La prononciation moderne permet d'utiliser des ressources audio vivantes (chansons, discours, liturgie orthodoxe) comme support de mémorisation.\cite{ref20}\item \textbf{La cohérence identitaire :} Prononcer le grec ancien \og à la moderne\fg{} (approche Reuchlinienne) réintègre les textes antiques dans le patrimoine vivant. Cela permet à l'élève de lire une inscription sur l'Acropole et de la prononcer de manière intelligible pour un Grec d'aujourd'hui, créant une gratification sociale immédiate.\cite{ref15}\item \textbf{L'argument musical :} L'intégration de la musique (même moderne ou inspirée de l'antique comme dans les jeux vidéo \textit{Assassin's Creed Odyssey}) nécessite une prononciation fluide qui respecte la mélodie de la langue telle qu'elle a évolué et survécu, plutôt qu'une reconstruction théorique souvent cacophonique.\cite{ref32}\end{itemize}

\subsection{Méthodologies actives et vivantes : études de cas}

La refondation cognitive ne reste pas théorique ; elle est mise en œuvre par des institutions
pionnières qui utilisent le grec moderne ou des méthodes de \og langues vivantes\fg{} pour enseigner
l'ancien. Ces modèles démontrent la viabilité de l'approche.



\subsection{L'institut polis (jérusalem) : la méthode rico et l'immersion totale}

L'Institut Polis représente la pointe avancée de cette pédagogie.\cite{ref29} Sous la direction de
Christophe Rico, linguiste agrégé, l'institut enseigne la Koine comme une langue vivante parlée.

\begin{itemize}\item \textbf{Le principe d'Immersion :} Les cours se déroulent intégralement en grec ancien. Cette immersion force le cerveau à traiter la langue comme un outil de communication et non comme un objet d'analyse, activant les mécanismes naturels d'acquisition (SLA \- Second Language Acquisition).\cite{ref37}\item \textbf{Living Sequential Expression (LSE) :} Rico a réactualisé une technique développée au XIXe siècle par le pédagogue François Gouin.\cite{ref37} Le LSE consiste à enseigner la langue par des séquences d'actions logiquement connectées (\og Je me lève, je marche vers la porte, j'ouvre la porte...\fg{}). Cette méthode s'appuie sur la séquentialité cognitive de l'action pour structurer la mémoire verbale.\item \textbf{L'usage pragmatique du moderne :} Bien que l'objectif soit le grec ancien, la méthode Polis n'hésite pas à emprunter au grec moderne pour nommer des réalités contemporaines nécessaires à la vie de classe (le téléphone, l'avion, etc.), marquant ces mots d'un astérisque, ou à adopter la prononciation historique pour maintenir le lien avec la culture vivante.\cite{ref28}\end{itemize}

\subsection{Elliniki agogi (athènes) : l'héritage vivant}

En Grèce, l'organisation \textit{Elliniki Agogi} propose une approche différente mais convergente,
centrée sur la continuité culturelle pour les enfants et les adultes.\cite{ref39}

\begin{itemize}\item \textbf{Philosophie \og Hellenizein\fg{} :} L'apprentissage est conçu comme une réappropriation identitaire. Le programme s'aligne sur le cursus scolaire moderne, utilisant la compétence native des élèves en grec moderne pour éclairer le grec ancien.\cite{ref42}\item \textbf{Méthode Expérientielle :} L'enseignement passe par le théâtre, le jeu, la récitation et le chant. Cette approche multisensorielle ancre le lexique ancien dans une expérience corporelle et émotionnelle positive, contrastant avec l'aridité scolaire traditionnelle.\cite{ref42}\end{itemize}Ces deux exemples montrent que le traitement du grec ancien comme une langue \og vivante\fg{} (ou en
continuité avec le vivant) permet de contourner les blocages cognitifs liés à l'abstraction
grammaticale.



\subsection{Le levier du vocabulaire scientifique : ancrage interne et utilité perçue}

Si le grec moderne est la ressource \og externe\fg{} (langue de la Grèce), le vocabulaire
scientifique est la ressource \og interne\fg{} déjà présente, souvent à l'état latent, dans l'esprit
de l'élève francophone. L'exploitation systématique de ce fonds lexical constitue le second pilier
de la motivation.



\subsection{Le grec, langue invisible de la modernité scientifique}

Il est établi que \textbf{80\% à 90\%} du vocabulaire scientifique et technique international
(anglais, français) dérive de racines gréco-latines.\cite{ref45} Ce \og grec scientifique\fg{} n'est
pas une relique, mais un système productif qui génère constamment des néologismes (ex:
\textit{biodégradable}, \textit{téléphone}, \textit{nanotechnologie}).

L'analyse étymologique transforme la nature de l'apprentissage lexical. Au lieu de mémoriser
arbitrairement que \textit{hydor} signifie \og eau\fg{}, l'élève connecte ce mot à un réseau
neuronal préexistant : \textit{hydraulique}, \textit{hydrogène},
\textit{déshydratation}.\cite{ref48}

\begin{itemize}\item \textbf{Effet de consolidation :} L'apprentissage du mot grec renforce la maîtrise du vocabulaire français et scientifique.\item \textbf{Effet de rétention :} Le lien sémantique profond (étymologie) est plus résistant à l'oubli que la traduction simple.\end{itemize}

\subsection{Contexte institutionnel et didactique (lca en france)}

En France, les programmes officiels (Eduscol) pour les Langues et Cultures de l'Antiquité (LCA) au
collège et au lycée ont intégré cette dimension comme un axe prioritaire.\cite{ref50}

\begin{itemize}\item \textbf{Interdisciplinarité :} Les instructions encouragent les projets croisés entre professeurs de LCA et professeurs de sciences (SVT, Physique-Chimie, Mathématiques). L'objectif est de montrer que le grec est la langue de la rationalité scientifique.\cite{ref50}\item \textbf{Légitimité pédagogique :} Les revues spécialisées comme les \textit{Cahiers Pédagogiques} soulignent que le travail sur l'étymologie et les racines est l'un des rares leviers capables de motiver des élèves en difficulté scolaire ou ceux qui se définissent comme \og scientifiques\fg{} et rejettent les humanités littéraires.\cite{ref54} Le grec devient un outil de décodage du monde, une compétence transversale valorisable.\end{itemize}

\subsection{Modélisation didactique : la "boucle étymologique"}

Pour la thèse, nous proposons de modéliser cette approche par la \og Boucle Étymologique\fg{}, qui
intègre ancien, moderne et scientifique :

\begin{table}[h!]\small\centering\begin{tabular}{| p{2cm} | p{3.5cm} | p{3.5cm} |}\hline\textbf{Étape} & \textbf{Processus} & \textbf{Exemple : Phon-} \\ \hline\textbf{1. Ancrage} & Partir du mot français courant & \textit{Smartphone}, \textit{Téléphone} \\ \hline\textbf{2. Analyse} & Isoler la racine et son sens & \textit{Phônè} = Voix, Son \\ \hline\textbf{3. Ancien} & Identifier dérivés en Grec Ancien & \textit{Phonéo} (parler) \\ \hline\textbf{4. Moderne} & Vérifier survivance en Grec Moderne & \textit{Tilefono}, \textit{Megafono} \\ \hline\textbf{5. Synthèse} & Évolution du concept & Voix humaine $\to$ transmission du son \\ \hline\end{tabular}\end{table}Cette démarche circulaire valide l'utilité du grec ancien non pas comme une fin en soi, mais comme
une clé de compréhension universelle.

La section III.4.\cite{ref2} de la thèse, en explorant le grec moderne comme ressource, met en
lumière une voie de refondation cognitive puissante. L'enseignement du grec ancien ne doit plus être
pensé comme l'exploration d'une nécropole linguistique, mais comme l'accès à une structure vivante
et résiliente.

Les preuves linguistiques de la continuité (Horrocks, Kavčič), combinées aux données
psycholinguistiques sur le seuil de lecture (Laufer \& Nation), démontrent que la segmentation
historique est un obstacle pédagogique artificiel. En s'appuyant sur le recouvrement lexical massif
(\~70-80\%) et en adoptant des méthodologies actives inspirées des langues vivantes (Polis),
l'enseignant peut réduire drastiquement la charge cognitive de ses élèves. Parallèlement,
l'exploitation du vocabulaire scientifique et la connexion avec la Grèce contemporaine offrent des
ancrages motivationnels indispensables dans le contexte éducatif actuel.

Le grec ancien, ainsi \og refondé\fg{}, redevient ce qu'il a toujours été : non pas une langue
morte, mais la racine vivante de la pensée occidentale et de la langue hellénique d'aujourd'hui.



\subsection{Tableau 2 : comparaison des approches de la prononciation}

26

\begin{table}[h!]\small\centering\begin{tabular}{| p{2cm} | p{3cm} | p{3cm} |}\hline\textbf{Caract.} & \textbf{Érasmienne} & \textbf{Historique/Moderne} \\ \hline\textbf{Origine} & Reconstruction théorique (XVIe s.) & Usage byzantin et moderne \\ \hline\textbf{Voyelles} & $\eta$=ê, $\iota$=i, $\upsilon$=u/y & Iotacisme ($\eta$, $\iota$, $\upsilon$, $\epsilon\iota$, $o\iota$ = i) \\ \hline\textbf{Consonnes} & B=b, D=d & B=v, D=ð \\ \hline\textbf{Avantage} & Correspondance graphie-phonie simple & Accès aux sources modernes \\ \hline\textbf{Inconvénient} & Rupture avec la Grèce actuelle & Homophonies nombreuses \\ \hline\end{tabular}\end{table}

\subsection{Tableau 3 : impact de la continuité lexicale sur l'apprentissage (données synthétiques kavčič \& laufer)}

\begin{table}[h!]\small\centering\begin{tabular}{| p{2cm} | p{2.2cm} | p{2cm} | p{2.3cm} |}\hline\textbf{Niveau} & \textbf{Voc. Ancien \%} & \textbf{À acquérir} & \textbf{Effort} \\ \hline\textbf{Débutant} & 0\% & 95\% & Maximal \\ \hline\textbf{Intermédiaire} & 30-40\% & 60\% & Moyen \\ \hline\textbf{Avancé} & 70-80\% & 15-25\% & Faible \\ \hline\end{tabular}\end{table}